\subsection*{Chapter 20: Making Inference Cheaper and Deploying It}

\paragraph{Exercise 20.1.} Qwen3-0.6B stores $2\cdot28\cdot8\cdot128\cdot2=114{,}688$ bytes
of KV per token, i.e.\ 235\,MB per 2048-token sequence. Its weights take 1.19\,GB, so the
two are equal at $B\approx1.19\cdot10^9/2.35\cdot10^8\approx5$. Beyond a handful of
concurrent sequences the KV cache dominates the bytes per step. Weight-only quantization
of a small model therefore barely helps a busy server; KV quantization, paging and
prefix sharing matter far more.

\paragraph{Exercise 20.2.} With $e$ uniform on $[-s/2,s/2]$,
$\E[e^2]=\frac1s\int_{-s/2}^{s/2}e^2\,de=s^2/12$. A full-range signal uniform on $[-A,A]$
has power $A^2/3$, and $b$ bits give $s=2A/2^b$, so the noise power is
$A^2/(3\cdot4^b)$. Hence $\mathrm{SQNR}=4^b$, or $10\log_{10}4^b\approx6.02\,b$\,dB.

\paragraph{Exercise 20.3.} Per row: $s=40/7\approx5.71$, and a value rounds to zero when
$|x|<s/2\approx2.86$, which holds for $\approx99.6\%$ of standard normal draws. The row is
destroyed. With $g=32$, only the outlier's group (31 normal values) suffers this. In the
other groups the scale is set by the group's own maximum ($\approx2.3$), so only
$|x|\lesssim0.17$ rounds to zero. A simulation gives $\approx34\%$ zeros overall, almost
all of them in the outlier's group.

\paragraph{Exercise 20.4.} (a) $4+(16+16)/64=4.5$ bits. (b) $4+8/32=4.25$ bits.
(c) $4+8/16=4.5$ bits, plus a negligible per-tensor fp32 scale.

\paragraph{Exercise 20.5.} For one row, the loss is
$\ell(\vect w')=(\vect w'-\vect w)^\top H(\vect w'-\vect w)$. Minimise $\ell$ subject to
$\vect e_i^\top(\vect w'-\vect w)=q_i-w_i$. The Lagrangian gives
$2H\boldsymbol\delta=\lambda\vect e_i$, so $\boldsymbol\delta=\frac{\lambda}{2}H^{-1}\vect e_i$,
and the constraint fixes $\lambda/2=(q_i-w_i)/[H^{-1}]_{ii}$. Therefore
$\boldsymbol\delta=-\frac{w_i-q_i}{[H^{-1}]_{ii}}H^{-1}_{:,i}$, whose $i$-th entry is
exactly $q_i-w_i$ and whose other entries are the stated update. Substituting back,
$\ell=\boldsymbol\delta^\top H\boldsymbol\delta=(w_i-q_i)^2/[H^{-1}]_{ii}$. GPTQ then
removes row and column $i$ from $H^{-1}$ (a Schur-complement update), which is what the
Cholesky formulation does implicitly.

\paragraph{Exercise 20.6.}
$(X\,\mathrm{diag}(\vect s)^{-1})(W\,\mathrm{diag}(\vect s))^\top
=X\,\mathrm{diag}(\vect s)^{-1}\mathrm{diag}(\vect s)W^\top=XW^\top$.
With $\alpha=\tfrac12$, $s_j=\sqrt{a_j/w_j}$ where $a_j=\max|X_{:,j}|$ and
$w_j=\max|W_{:,j}|$. Then $\max|X'_{:,j}|=a_j/s_j=\sqrt{a_jw_j}=w_js_j=\max|W'_{:,j}|$.
To fold the division, note that $X$ is the output of RMSNorm with per-channel gain
$\gamma_j$; replace $\gamma_j$ by $\gamma_j/s_j$ offline. (When several linears consume
the same normalised input, e.g.\ Q, K and V, they must share $\vect s$.)

\paragraph{Exercise 20.7.} Let $x_1$ be emitted at position 1 with probability $p(x_1)$,
by the proposition. Conditioned on the first draft having been accepted with value $x_1$,
the second draft is sampled from $q(\cdot\mid x_1)$ and verified against
$p(\cdot\mid x_1)$, because the target's second output row is computed with $x_1$ in its
context. Applying the proposition to this conditional pair gives
$\Pr[x_2\mid x_1]=p(x_2\mid x_1)$. If the first draft is rejected, the round ends and
position 2 is produced in the next round, again conditioned on the emitted $x_1$.
Either way the joint law is $p(x_1)p(x_2\mid x_1)$. Induction extends this to whole
sequences.

\paragraph{Exercise 20.8.} Maximising
$S(\gamma)=(1-\alpha^{\gamma+1})/((1-\alpha)(0.05\gamma+1))$ over integers gives
$\gamma^{*}=4$ ($S\approx1.92$) for $\alpha=0.6$, $\gamma^{*}=8$ ($S\approx3.09$) for
$\alpha=0.8$, and $\gamma^{*}=13$ ($S\approx4.67$) for $\alpha=0.9$. Going from $\gamma$
to $\gamma+1$ helps iff
$S(\gamma+1)>S(\gamma)$, i.e.\
$\alpha^{\gamma+1}(1-\alpha)(c\gamma+1)>c\,(1-\alpha^{\gamma+1})$: the expected marginal
token $\alpha^{\gamma+1}$ must outweigh the marginal drafting cost relative to the
current expected yield.

\paragraph{Exercise 20.9.} The step reads $16.4$\,GB of weights plus
$128\cdot2048\cdot147.5$\,KB$\approx38.7$\,GB of KV: $t_{\text{mem}}\approx16.4$\,ms. With
$\gamma+1=5$, the step processes 640 tokens: $\approx2\cdot8.19\cdot10^9\cdot640\approx
1.05\cdot10^{13}$ matmul FLOPs, plus $\approx0.8\cdot10^{12}$ for attention, so
$t_{\text{comp}}\approx11.4$\,ms at peak. That is still memory-bound on paper, but at a
realistic 50--60\% of peak the compute time reaches $\approx20$\,ms, which exceeds the
memory time. Every rejected draft then costs real time. Choose $\gamma\le2$ (or adapt
$\gamma$ to load), and turn speculation off entirely at higher batch sizes.

\paragraph{Exercise 20.10.} The message size is $m=5120\cdot2=10$\,KB per sequence, which
is negligible on NVLink ($m/\beta\approx0.02\,\mu$s). If the per-collective latency is
constant ($a=8\,\mu$s, one-shot all-reduce),
$t(P)=19.6/P+2\cdot64\cdot0.008$\,ms$=19.6/P+1.02$\,ms. This decreases monotonically:
10.8\,ms at $P=2$, 5.9\,ms at $P=4$, 3.5\,ms at $P=8$. So use the whole NVLink domain. If
the latency instead grows like a ring all-reduce, $a(P)=2(P-1)\alpha$ with
$\alpha=2\,\mu$s, then $t=10.3$, $6.4$ and $6.0$\,ms for $P=2,4,8$: diminishing returns,
and a two-node $P=16$ would be slower. Collective latency, not bandwidth, sets the limit
on TP for decode.

\paragraph{Exercise 20.11.} In a closed loop a new request is issued only when an old one
completes, so the number in the system is capped at 64 and the offered load automatically
drops when the server slows down. No queue can build up, and latency stays bounded by
construction. Real traffic is open-loop: Poisson bursts arrive regardless of server
state. During a burst the queue and the KV cache fill up, and preemption and queueing
delay grow superlinearly near capacity (Little's law plus queueing theory). Closed-loop
tools also suffer from ``coordinated omission'': the slow periods generate fewer samples,
so they are under-represented in the percentiles.

\paragraph{Exercise 20.12.} Assume every message is 150 tokens. The prompt at turn $k$ is
about $300(k-1)+150$ tokens, so without reuse the total prefill is
$\sum_{k=1}^{20}(300k-150)=60{,}000$ tokens. With perfect reuse each turn prefills only
the new user message plus a few template tokens (and the one generated token that was
never fed back): about $20\cdot150=3{,}000$ tokens, a $20\times$ reduction. With the
official Qwen3 template, pitfall (1) strikes first: the earlier assistant turns'
\texttt{<think>} blocks are stripped when the history is re-serialised, so the common
prefix ends at the start of the most recent assistant reply. Fixes: store in the history
exactly what the template will emit (the reply without its think block) and accept
re-prefilling that shorter reply once, or keep a cache keyed on the template's output.
Engine-level prefix caching (radix trees) handles this automatically, at block
granularity.
